% TOP_PROBLEM: {"id": 30007232, "title": "A size-independent spectral gap above the spherical Laughlin line", "category_id": 16, "category": {"id": 16, "name": "physics", "display_name": "Mathematical Physics", "description": "Problems at the intersection of mathematics and physics.", "slug": "physics", "order_index": 16, "created_at": "2026-07-31T15:26:25.671Z"}, "status": "open"}
% FIRST_POSTED: 2026-10-06
% ATTEMPT_STATUS: solved
% OPENAI_CLAIM: {"schema_version":1,"source_repo":"https://github.com/openai/math","source_commit":"adc7f1241b42e322a6451854ab7e4b4c146bf78a","release_date":"2026-10-06","match":"full","resolution":"proved","independently_reviewed":false,"second_pass":"agreed","families":[{"family":"269","title":"Uniform Laughlin gap and stability under bounded scalar disorder","subject":"Mathematical physics","manuscripts":[{"title":"A Fock-space inequality and the Laughlin spectral gap","pdf_path":"preprints/A-Fock-space-inequality-and-the-Laughlin-spectral-gap-September-24-2026/A-Fock-space-inequality-and-the-Laughlin-spectral-gap-September-24-2026.pdf","readme_path":"preprints/A-Fock-space-inequality-and-the-Laughlin-spectral-gap-September-24-2026/README.md","date":"2026-09-24","theorem_ref":"thm:main (sections/01-introduction.tex:60)","statement_sources":[{"path":"preprints/A-Fock-space-inequality-and-the-Laughlin-spectral-gap-September-24-2026/build/sections/01-introduction.tex","line":60,"label":"thm:main"}]}],"lean":{"doc_path":"lean/docs/269.md","comparators":["lean/ComparatorChallenges/Laughlin.json","lean/ComparatorChallenges/Laughlin.lean","lean/ComparatorChallenges/LaughlinFock.json","lean/ComparatorChallenges/LaughlinFock.lean","lean/ComparatorChallenges/LaughlinGap.json","lean/ComparatorChallenges/LaughlinGap.lean","lean/ComparatorChallenges/LaughlinPlanar.json","lean/ComparatorChallenges/LaughlinPlanar.lean"],"declarations":["OAI.Laughlin.mainTarget_proved"],"covers_main_theorem":true},"reasoning_trace":null}]}
\documentclass[11pt]{article}
\usepackage[margin=25mm]{geometry}
\usepackage{fontspec}
\setmainfont{Latin Modern Roman}
\usepackage{amsmath,amssymb,mathtools}
\usepackage{unicode-math}
\setmathfont{Latin Modern Math}
\usepackage{enumitem}
\usepackage{xurl,hyperref,bookmark}
\hypersetup{colorlinks=true,urlcolor=blue}
\setcounter{secnumdepth}{0}
\setlength{\parindent}{0pt}
\setlength{\parskip}{5pt}
\setlength{\emergencystretch}{3em}
\begin{document}
\section*{A size-independent spectral gap above the spherical Laughlin line}

\subsection{OpenAI's claim}
OpenAI presents a proof in the selected manuscript. The spherical fermionic Laughlin $V_1$ model is claimed uniformly gapped at filling $1/3$, with a stronger lower bound on the square of the interaction across the entire lowest-level Fock space.

\subsection{Mathematical statement of the claimed result}
\paragraph{Result 1.} For the lowest-Landau-level pseudopotential \(H_{N,Q}=\sum_{i<j}P_{ij}^{(1)}\) on \(\bigwedge^NU_Q\), with pair total-spin projector \(Q-1\) and coefficient one, the source claims an integer \(N_0\) such that
\[
 H_{N,3(N-1)}\ge\frac1{25}(I-P_{\mathrm L,N})\qquad(N\ge N_0).
\]
Here \(P_{\mathrm L,N}\) projects onto the line spanned by \(\prod_{i<j}(u_iv_j-u_jv_i)^3\). Thus the spherical Laughlin line is a unique zero-energy ground space with a size-independent positive excitation gap.

\subsection{Outline of the argument}
Exterior-algebra identities expand the interaction square into spin coefficients. Uniform asymptotic control of these coefficients bounds cross terms against the original interaction, giving a Fock-space inequality independent of particle number and hence the Laughlin gap.

\subsection{Scope relative to this problem}
The spherical fermionic Laughlin $V_1$ model is claimed uniformly gapped at filling $1/3$, with a stronger lower bound on the square of the interaction across the entire lowest-level Fock space. The question retains these same parameters and conventions. Within these assumptions, this is a claimed affirmative answer to the stated question.

\subsection{Formal verification}
For family 269, the documented entries give unperturbed spherical Laughlin and stronger Fock-space gap inequalities. Stability against scalar one-body disorder and uniqueness of the perturbed ground state are not included. The listed Comparator files are \texttt{Laughlin.lean}, \texttt{LaughlinFock.lean}, \texttt{LaughlinGap.lean}, \texttt{LaughlinPlanar.lean}. Named declarations include \texttt{OAI.Laughlin.mainTarget\_proved}. This site has not executed the Lean code. The scope descriptions do not certify the remaining manuscript arguments.

\subsection{Status caveat}
OpenAI attributes these manuscripts to work by an internal model. Its release cautions that arguments without formalization can contain mistakes. This catalogue has not independently checked the proofs; the labels report the mathematical claims attributed to OpenAI.

\subsection{Sources}
\begin{itemize}
\item OpenAI, \emph{A Fock-space inequality and the Laughlin spectral gap}, 2026-09-24; citation key \texttt{OAI:A-Fock-space-inequality-and-the-Laughlin-spectral-gap-September-24-2026}. \href{https://github.com/openai/math/blob/adc7f1241b42e322a6451854ab7e4b4c146bf78a/preprints/A-Fock-space-inequality-and-the-Laughlin-spectral-gap-September-24-2026/A-Fock-space-inequality-and-the-Laughlin-spectral-gap-September-24-2026.pdf}{PDF (GitHub)}; \href{https://github.com/openai/math/raw/adc7f1241b42e322a6451854ab7e4b4c146bf78a/preprints/A-Fock-space-inequality-and-the-Laughlin-spectral-gap-September-24-2026/A-Fock-space-inequality-and-the-Laughlin-spectral-gap-September-24-2026.pdf}{PDF (direct)}; \href{https://github.com/openai/math/blob/adc7f1241b42e322a6451854ab7e4b4c146bf78a/preprints/A-Fock-space-inequality-and-the-Laughlin-spectral-gap-September-24-2026/README.md}{Manuscript page and citation}; \href{https://github.com/openai/math/blob/main/preprints/A-Fock-space-inequality-and-the-Laughlin-spectral-gap-September-24-2026/A-Fock-space-inequality-and-the-Laughlin-spectral-gap-September-24-2026.pdf}{Latest version}; Statement source: \href{https://github.com/openai/math/blob/adc7f1241b42e322a6451854ab7e4b4c146bf78a/preprints/A-Fock-space-inequality-and-the-Laughlin-spectral-gap-September-24-2026/build/sections/01-introduction.tex\#L60}{thm:main}.
\item Family 269: \href{https://github.com/openai/math/blob/adc7f1241b42e322a6451854ab7e4b4c146bf78a/lean/docs/269.md}{Lean scope}.
\item \href{https://github.com/openai/math/blob/adc7f1241b42e322a6451854ab7e4b4c146bf78a/lean/ComparatorChallenges/Laughlin.json}{Comparator Laughlin.json}.
\item \href{https://github.com/openai/math/blob/adc7f1241b42e322a6451854ab7e4b4c146bf78a/lean/ComparatorChallenges/Laughlin.lean}{Comparator Laughlin.lean}.
\item \href{https://github.com/openai/math/blob/adc7f1241b42e322a6451854ab7e4b4c146bf78a/lean/ComparatorChallenges/LaughlinFock.json}{Comparator LaughlinFock.json}.
\item \href{https://github.com/openai/math/blob/adc7f1241b42e322a6451854ab7e4b4c146bf78a/lean/ComparatorChallenges/LaughlinFock.lean}{Comparator LaughlinFock.lean}.
\item \href{https://github.com/openai/math/blob/adc7f1241b42e322a6451854ab7e4b4c146bf78a/lean/ComparatorChallenges/LaughlinGap.json}{Comparator LaughlinGap.json}.
\item \href{https://github.com/openai/math/blob/adc7f1241b42e322a6451854ab7e4b4c146bf78a/lean/ComparatorChallenges/LaughlinGap.lean}{Comparator LaughlinGap.lean}.
\item \href{https://github.com/openai/math/blob/adc7f1241b42e322a6451854ab7e4b4c146bf78a/lean/ComparatorChallenges/LaughlinPlanar.json}{Comparator LaughlinPlanar.json}.
\item \href{https://github.com/openai/math/blob/adc7f1241b42e322a6451854ab7e4b4c146bf78a/lean/ComparatorChallenges/LaughlinPlanar.lean}{Comparator LaughlinPlanar.lean}.
\item \href{https://github.com/openai/math/blob/adc7f1241b42e322a6451854ab7e4b4c146bf78a/CONTENTS.md}{OpenAI catalogue entry}.
\end{itemize}
\end{document}
